% Generated draft: requires mathematical and source review.
\documentclass[11pt]{article}
\usepackage[margin=25mm]{geometry}
\usepackage{fontspec}
\setmainfont{Latin Modern Roman}
\usepackage{amsmath,amssymb,mathtools,unicode-math}
\setmathfont{Latin Modern Math}
\usepackage{enumitem,needspace,xurl,hyperref}
\setlength{\parindent}{0pt}
\setlength{\parskip}{5pt}
\setlength{\emergencystretch}{3em}
\newcommand{\N}{\mathbb{N}}
\newcommand{\Z}{\mathbb{Z}}
\newcommand{\Q}{\mathbb{Q}}
\newcommand{\R}{\mathbb{R}}
\newcommand{\C}{\mathbb{C}}
\newcommand{\F}{\mathbb{F}}
\newcommand{\A}{\mathbb{A}}
\newcommand{\PP}{\mathbb{P}}
\newcommand{\E}{\mathbb{E}}
\newcommand{\eps}{\varepsilon}
\newcommand{\dd}{\,\mathrm{d}}
\begin{document}
\section*{UnsolvedMath ID: 5300037}
\textbf{Source problem code:} catalogue entry
\subsection{Definitions and mathematical statement}
For $c\in\C$, let $f_c(z)=z^2+c$, let
\[
 K(f_c)=\{z\in\C:(f_c^n(z))_{n\ge0}\text{ is bounded}\},
 \qquad J(f_c)=\partial K(f_c).
\]
A compact set is locally connected when each point has a basis of connected relative neighborhoods.

Quadratic-like renormalization refers to a proper degree-two restriction $f_c^p:U\to V$ between topological disks with $\overline U\subset V$, capturing the critical dynamics and having connected filled Julia set. Infinite renormalizability requires renormalizations at arbitrarily deep nested levels. The source's additional condition ``of bounded type'' is not defined in the record. It commonly involves bounded successive renormalization combinatorics, such as bounded relative return periods, but the exact convention and any primitive/satellite scope must be checked rather than silently selected. No a priori bound is added as an extra hypothesis.

The question is whether
\[
 \forall c\in\C,\qquad
 \bigl(f_c\text{ is infinitely renormalizable of bounded type}\bigr)
 \ \Longrightarrow\ J(f_c)\text{ is locally connected}.
\]
It explicitly includes the special question whether $J(f_{c_F})$ is locally connected for the quadratic Feigenbaum map, with $c_F$ the usual real period-doubling accumulation parameter. The record supplies no alternative normalization or numerical parameter value.

The object in the conclusion is the individual dynamical Julia set. The research summary uses ``MLC'' while discussing this result, but local connectivity of the Mandelbrot parameter set is a different assertion and is not substituted for the displayed conclusion.
\subsection{Short English statement}
Are Julia sets locally connected for all infinitely renormalizable quadratic maps of bounded type, in particular the quadratic Feigenbaum map? The supplied research summary answers yes and labels the record solved in the literature. It attributes the result to a priori bounds and renormalization methods involving Kahn, Lyubich, and Kahn--Lyubich, but gives no exact theorem statements or complete references.

The summary contrasts unbounded-type renormalization, including growing satellite combinatorics, as a subtler setting and refers to AMR-052-0012; it does not extend the claimed bounded-type result to that setting. Provenance says the 1992 Stony Brook source presents the problem as open and that no comprehensive modern resolution check was available. The solved assessment is therefore retained with those source-verification limits.
\subsection{Sources}
\begin{itemize}
\item Local UnsolvedMath registry, permanent ID 5300037, source code catalogue entry, ``Local connectivity for bounded-type renormalization,'' proposed by John Milnor, 1992. Dataset curation and metadata: Ulam AI and UnsolvedMath Contributors, CC BY 4.0.
\item Stony Brook, \emph{Open Problems in Dynamical Systems}, 1992, Milnor local P12, \url{https://www.math.stonybrook.edu/open-problems-dynamical-systems}; source-TeX and pdf-text extraction, accessed 2026-07-29.
\item Supplied research-summary attributions: Kahn; Lyubich; and Kahn--Lyubich, for a priori bounds, renormalization methods, and local connectivity. No titles, dates, exact hypotheses, or publication locators are supplied.
\item The summary refers to the separate registry source code AMR-052-0012 for unbounded-type questions; no permanent numeric ID is inferred from that code.
\item Source provenance marks status evidence NEEDS\_REVIEW and rights NEEDS\_REVIEW: public preprint access and provenance were verified, but redistribution permission and a comprehensive modern resolution check were not established. This qualification is separate from the dataset licence.
\end{itemize}
\end{document}
